% =============================================================================
\chapter{One paper by Friday: the merge}
\label{ch:merge}

The abstract and the author list lock on Friday 18 September, AoE; the paper is due on Friday 25
September, AoE. Your draft has the object, the geometry, the theory and the algorithm. Our side has an
evaluation that already ran, on real structure, with the audit opened after the rows were frozen. This
chapter maps one onto the other: what your open items correspond to, what the measurements add, where the
two drafts disagree, three ways to merge them, and the order of work for the next nine days.

\begin{ideia}[The merge in one sentence]
Keep your object, geometry and theory as they stand; replace the evaluation section with the protocol that
already ran; move the title and abstract to the pricing frame of Rosenbaum's $\Gamma$; and cut to nine
pages.
\end{ideia}

\begin{antes}
\antesq[pre-registro]{Your draft owes two experiments in \texttt{\textbackslash todo} markers. Which of
them has already run, and with which endpoint?}
\antesq[unidade-afirmacao]{Your commented-out caveat says $\rclaim = 1$ on all full-coverage rows by
construction. What removed that construction?}
\end{antes}

\section{Your open items are the runs of 11 to 13 September}

Your draft carries three \texttt{\textbackslash todo} markers. The first asks for a single end-to-end
speedup figure; nothing on our side answers it. The other two correspond to the two measurements of
Chapters~\ref{ch:frame} to~\ref{ch:table4}, the second one almost word for word.

\textbf{\S6.2.} «The ranking evaluation of Section~6.1 has not yet been run on these 831 real instances;
survival and the paired cross-arm design exist only on synthetic structure.» What ran on the committed real
rows is the corrupted-knowledge calibration of Chapter~\ref{ch:frame}, the nearest thing to what the
marker asks. It flips claims of the recovering set at four rates and two coverages on the same real rows,
and asks of every instrument whether its certificate held, scored in the four buckets of
Chapter~\ref{ch:vocabulary} rather than by $\tau_b$ against the survival area. The $\tau_b$ design itself
has not been run on the real rows.

\textbf{\S6.3.} «The elicitation arm: knowledge collected by asking for a causal order over each chain
component rather than for pairwise edge directions, from which $\rclaim$, $\rval$ and $\varphi_1$ are
measured on the same frozen instance frame, with a matched random control and an audit opened only after
the rows are frozen.» That is the ledger, exactly as specified.

\begin{medido}[The two markers, closed: \texttt{results/\allowbreak stage0/\allowbreak } and \texttt{results/\allowbreak ledger/\allowbreak }]
\textbf{Calibration} (\texttt{results/\allowbreak stage0/\allowbreak calibration\_summary.json}): 16,496 draws, 8,911 scored,
24 networks, flip rates 0, 0.1, 0.25 and 0.5 at coverages 1.0 and 0.5, five repetitions per nonzero rate,
seeds derived by SHA-256. \texttt{results/\allowbreak stage0/\allowbreak STAGE0.md} still quotes 8,553 scored rows from the run before the
seeding fix; the JSON is the rerun.
\textbf{Elicitation and ledger}: the questionnaire covers 120 chain components, 369 open pairs and 540
rendered items, frozen and hashed before any supplier saw it (\texttt{docs/\allowbreak PROTOCOL\_RUN.md});
\texttt{results/\allowbreak ledger/\allowbreak rows.csv} holds 8,308 rows over 14 knowledge conditions, with the chance control
\texttt{D\_RAND} matched on the supplier's edges and $\rclaim$, $\rhop$, $\varphi_1$ on 528 frame rows per
supplier condition (\texttt{results/\allowbreak ledger/\allowbreak TABLE4.md}).
\end{medido}

Your commented-out caveat in \S6.3 is also right, and the protocol removed its cause. \texttt{docs/\allowbreak EVALUATION\_PLAN.md}
confirms that the admissibility gate makes $\rclaim = 1$ on 540 of 540 full-coverage rows. The frame of
Chapter~\ref{ch:frame} is drawn from the CPDAG alone, 80,132 candidate pairs over 33 networks, and contains
every one of the 543 committed pairs, so $\rclaim$ is now measured rather than fixed by construction.

\prompt[conceito=pre-registro,dificuldade=0.7]{Map each of your two experimental \texttt{\textbackslash todo}
markers to the artefact that closes it, and say how the first one differs from what you asked for.}%
{The \S6.3 elicitation arm is the ledger: a causal-order questionnaire over 120 components (369 open
pairs, 540 items) frozen and hashed first, 8,308 rows over 14 conditions in \texttt{results/ledger/}, a
matched chance control and an audit opened after freezing. The \S6.2 real-structure evaluation is the
corrupted-knowledge calibration on the committed rows (8,911 scored of 16,496 draws over 24 networks, in
\texttt{calibration\_summary.json} under \texttt{results/stage0/}), which scores instruments in the four buckets rather than by
$\tau_b$ against the survival area; the $\tau_b$ design itself has not run on real rows.}

% -----------------------------------------------------------------------------
\begin{antes}
\antesq[dois-paineis]{Your survival table and our bucket audit measure the same object. May they share a
table?}
\antesq[screen-leave-k-out]{Which baseline would a referee add to your $\tau_b$ table first?}
\end{antes}

\section{What the measurements add, and where the drafts conflict}

Eight things exist on our side and not in your draft. Each has its chapter and its file.

\begin{enumerate}[nosep, leftmargin=*]
\item \textbf{The knowledge interval and the null radius} (Chapter~\ref{ch:gamma}). With one wrong claim
  the stated report covers 0.597 and the interval at one claim covers 0.980 at an $\Lnf$ of 0.851 (Panel~A,
  $n = 1000$, \texttt{results/\allowbreak interval/\allowbreak TABLE1\_L95.md}).
\item \textbf{The $\Gamma$ framing}: price what cannot be tested, with the interval as the headline
  (Chapter~\ref{ch:gamma}).
\item \textbf{The unit, measured}: over 3,014 certifiable rows (true CPDAG) the hop radius over-certifies on
  39, the free rule on 15, the claim certificate on none (\texttt{results/\allowbreak ledger/\allowbreak TABLE4.md}).
\item \textbf{The stop rule against the free screen} (Chapter~\ref{ch:m1b}): refusal 0.910 on 134
  informative committed rows with correct knowledge (Panel A); different decisions from the one-claim
  screen on 2 of 162 rows at one wrong claim (both panels); and a coverage difference of
  $+0.000\ [+0.000, +0.000]$ against the pair screen at two wrong claims (Panel A, $n = 1000$) (\texttt{results/\allowbreak interval/\allowbreak TABLE1\_L95.md}).
\item \textbf{The estimated class} (Chapter~\ref{ch:pest}): the truth extends the analyst's graph on 0 of
  244 certifiable rows, and both certificates are dangerous on 14 of 38 (\texttt{results/\allowbreak pest/\allowbreak TABLE7.md}).
\item \textbf{The audit oracle}: the back-door oracle rejected 4 of 1,169 committed sets that the GAC
  accepts, so the audit now uses \texttt{src/\allowbreak bkrobust/\allowbreak benchmarks/\allowbreak audit.py}, guarded by
  \texttt{tests/\allowbreak benchmarks/\allowbreak test\_audit\_oracle.py} (\texttt{results/\allowbreak ledger/\allowbreak PREREGISTRATION.md}, A.0b).
\item \textbf{The suppliers as error processes} (Chapter~\ref{ch:suppliers}).
\item \textbf{The monotonicity audit}: 5,000 sampled pairs over 469 rows, 0 violations, a one-sided 95\,\%
  upper bound of 0.06\,\% (\texttt{results/\allowbreak ledger/\allowbreak monotonicity\_audit.json}): the implemented failure
  predicate is upward-closed, which is your Theorem~1 as run.
\end{enumerate}

Four conflicts follow, and a merged paper has to settle each of them in its text.

\textbf{Framing.} Your paper is a diagnostic radius; ours is a sensitivity analysis in the sense of
Rosenbaum's $\Gamma$, with the interval leading. Your abstract already says the radius «prices the
structural risk», which is the pricing move; what it prices is validity alone.

\textbf{The unit.} Your headline counts covering hops and treats claims as a refinement. The ledger says
the hop count over-certifies and the claim count does not, and on your own running example the two give
3 and~2 (Chapter~\ref{ch:vocabulary}).

\textbf{A missing baseline.} Your $\tau_b$ table has no leave-$k$-out screen, and M1b showed that at matched
depth the stop rule built on the radius takes the screen's decisions, except where a retraction changes
the set while leaving it valid.

\textbf{Two designs on one object.} Survival ranking on synthetic structure and the bucket audit on real
structure answer different questions on related populations. They go in the paper as two panels with
their populations named, never pooled, which is your own \S5.4 rule.

\prompt[conceito=screen-leave-k-out,dificuldade=0.7]{What must be added to your $\tau_b$ table, and what does
M1b say it will show at matched depth and against the depth-one screen?}%
{A leave-$k$-out screen: withdraw one (or two) claims, re-close, and flag if the committed result changes.
At matched depth the stop rule and the screen take the same decisions except where a retraction changes
the set while leaving it valid: they differ on 2 of 162 rows at one wrong claim (both panels) and on 5 of
754 queries at depth two, and the coverage difference against the pair screen at two wrong claims is
$+0.000\ [+0.000, +0.000]$ (Panel A, $n = 1000$).
Against the one-claim screen at two wrong claims, the two-claim stop rule covers $+0.256\ [+0.149, +0.346]$
more at $n = 1000$ ($+0.305$ in the population row). The licensed reading is that the radius tells the
analyst which depth was needed, not that it beats the screen.}

% -----------------------------------------------------------------------------
\begin{antes}
\antesq[decisoes-abertas]{Which sections of your draft survive a merge almost untouched?}
\antesq[deriva-enquadramento]{Can the title still change after 18 September?}
\end{antes}

\section{Three ways to merge, and the one I recommend}

\begin{figure}[!htbp]
\centering
\begin{tikzpicture}[font=\scriptsize\sffamily,
  dr/.style={draw=cinzaescuro, fill=white, rounded corners=2pt, inner sep=2.5pt, text width=3.55cm,
             align=left, minimum height=0.55cm},
  mg/.style={draw=azulesc!70, fill=azulesc!5, rounded corners=2pt, inner sep=2.5pt, text width=5.3cm,
             align=left, minimum height=0.55cm},
  ours/.style={draw=verdesc, fill=verdesc!10, rounded corners=2pt, inner sep=3pt, text width=2.9cm,
             align=left},
  keep/.style={-{Stealth[length=1.5mm]}, cinzamedio, thick},
  rew/.style={-{Stealth[length=1.5mm]}, laranjaesc, thick},
  rep/.style={-{Stealth[length=1.5mm]}, vermesc, thick},
  mov/.style={-{Stealth[length=1.5mm]}, cinzamedio, thick, densely dashed}]
\node[font=\small\sffamily\bfseries] at (0,0.75) {your draft, 11 pp};
\node[font=\small\sffamily\bfseries] at (6.75,0.75) {merged paper, option A, 9 pp};
\node[dr] (d0) at (0,0) {title, abstract};
\node[dr] (d1) at (0,-0.8) {\S1 introduction};
\node[dr] (d2) at (0,-1.6) {\S2 preliminaries};
\node[dr] (d3) at (0,-2.4) {\S3 related work};
\node[dr] (d4) at (0,-3.2) {\S4 distance, radius};
\node[dr] (d5) at (0,-4.0) {\S5 tractability, protocol};
\node[dr] (d6) at (0,-4.8) {\S6 simulation study};
\node[dr] (d7) at (0,-5.6) {\S7--8 discussion and conclusion};
\node[dr] (d9) at (0,-6.4) {appendices A--E};
\node[mg] (m0) at (6.75,0) {title and abstract: the pricing frame};
\node[mg] (m1) at (6.75,-0.8) {\S1 + clinician sentence on page one};
\node[mg] (m2) at (6.75,-1.6) {\S2, criterion named and attributed};
\node[mg] (m3) at (6.75,-2.4) {\S3 + breakdown frontiers, correction sets, adjustment distances, tiered knowledge};
\node[mg] (m4) at (6.75,-3.2) {\S4 + claim unit, Prop.~9 paragraph};
\node[mg] (m5) at (6.75,-4.0) {\S5, conditions stated, shortened};
\node[mg] (m6) at (6.75,-4.8) {\S6 executed protocol: frame, ledger, interval, screen, estimated class};
\node[mg] (m7) at (6.75,-5.6) {\S7 limits paragraph uncommented};
\node[mg] (m9) at (6.75,-6.4) {appendix: proofs, conjecture stated; survival $\tau_b$ panel with a screen column};
\node[ours] (o) at (12.35,-4.0) {our measurements,\\ Chapters~\ref{ch:frame} to~\ref{ch:tables}};
\draw[rew] (d0.east) -- (m0.west);
\draw[rew] (d1.east) -- (m1.west);
\draw[keep] (d2.east) -- (m2.west);
\draw[rew] (d3.east) -- (m3.west);
\draw[keep] (d4.east) -- (m4.west);
\draw[keep] (d5.east) -- (m5.west);
\draw[rep] (d6.east) -- node[above, font=\tiny\sffamily, text=vermesc]{replaced} (m6.west);
\draw[mov] (d6.east) to[out=-10, in=175] (m9.west);
\draw[rew] (d7.east) -- (m7.west);
\draw[keep] (d9.east) -- (m9.west);
\draw[-{Stealth[length=1.5mm]}, verdesc, thick] (o.west) -- (m6.east);
\draw[-{Stealth[length=1.5mm]}, verdesc, thick] (o.north) |- (m0.east);
\node[font=\tiny\sffamily, cinzamedio, anchor=west] at (-1.8,-7.25) {grey: kept \quad orange: rewritten \quad red: replaced \quad dashed: moved to the appendix};
\end{tikzpicture}
\caption{Under option A almost everything you wrote survives; the evaluation section is the one part that
is replaced, and your synthetic survival table moves to the appendix as its own panel.}
\label{fig:merge-map}
\end{figure}

There are three ways to merge, and \figref{fig:merge-map} shows where each part of your draft goes under
the first.

\textbf{A. One paper, your skeleton, our measurements.} Your \S\S2 to~5 stay almost as they are, with
the fixes of Chapter~\ref{ch:review}; \S6 is replaced by the executed protocol (frame, suppliers, ledger,
interval and null radius, the screen, the estimated class); title and abstract move to the pricing frame;
the main text is cut from eleven pages to nine. Cost: one new section, an abstract, a two-page cut.

\textbf{B. One paper, my outline, your theory as its object section.} The $\Gamma$ framing leads, your
geometry and tractability become \S2 with proofs in the appendix. Cost: more rewriting of your prose, and
the tractability results lose prominence.

\textbf{C. Two papers.} Yours at ICLR, ours elsewhere. Cost: both use one corpus and one definition, each is
weaker alone, and the elicitation arm is what puts yours on real structure.

I recommend~A. It keeps what you wrote, closes both experimental markers with measured results, and is the
only option that fits two days before the lock.

Nine decisions for the text were listed in my outline of 12--13 September, and three of them change the
abstract. \textbf{D1}, the title: name the analysis, not the radius. \textbf{D2}, the headline value: lead
with the knowledge interval and report $\rnul$ beside it as the number of claims the conclusion rests on,
named. \textbf{D9}, the framing after M1b: accept the measurement framing declared on 7 September for
exactly this tie. These are recommendations for the two of us to decide, not decisions taken.

\prompt[conceito=decisoes-abertas,dificuldade=0.7]{Under option A, which parts of your draft are kept, which
are rewritten, and which one is replaced by what?}%
{Kept: \S2 preliminaries, \S4 distance and radius, \S5 tractability, the appendices, each with the fixes of
the review (criterion named, Proposition 9 paragraph and claim unit in \S4, conditions stated in \S5, the
conjecture stated). Rewritten: title and abstract (pricing frame), \S1 (clinician sentence on page one),
\S3 (missing citations), \S7 (limits paragraph uncommented). Replaced: \S6, by the executed protocol
(frame, ledger, interval and null radius, screen, estimated class); the synthetic $\tau_b$ table moves to
the appendix as its own panel with a screen column.}

% -----------------------------------------------------------------------------
\begin{antes}
\antesq[frases-proibidas]{Name one sentence the paper must not contain, even though a result makes it
tempting.}
\antesq[pre-registro]{After 18 September, what about the author list can still change?}
\end{antes}

\section{Before Friday, and before the 25th}

\begin{fonte}[ICLR 2027 Author Guidelines, read on 16 September]
«Please note that no authors can be added or removed after the abstract submission deadline (author order
can be changed after the abstract submission deadline up to the full paper deadline).» «Abstracts submitted
by the abstract submission deadline must be genuine; placeholder or duplicate abstracts will be removed.»
«Please coordinate among the authors and make sure that everyone has an OpenReview profile before the
abstract submission deadline.» «all submissions must have at least one author who is registered to review
at least 3 papers». Titles «can be edited up to the paper submission deadline». «At the time of submission,
the main text should be 9 pages or fewer.» «Required: AI use statement.»
\end{fonte}

So the author list, and with it everyone's OpenReview profile, must be settled by 18 September; only its
order can move afterwards. In order, before Friday 18 September, AoE:

\begin{enumerate}[nosep, leftmargin=*]
\item Settle the author list and check every author's OpenReview profile.
\item Decide the framing (D9) and the headline value (D2), and choose the merge option.
\item Narrow the novelty sentence to what was searched. The abstract's «nothing currently tells the analyst
  how much of the conclusion rests on them» is too broad: Taeb, Guo and Henckel's \S6 already reports how
  large a misspecification ball around the analyst's graph each adjustment set tolerates (one example, in
  ADMGs, by adding edges to the specified graph). Chapter~\ref{ch:tables} lists the novelty search still owed
  before the abstract.
\item Write the abstract from the merged claims, one sentence each for the interval, the radius in claims,
  the tie with the screen and the class assumption, and check it against the list below.
\item Submit title, abstract, authors and subject areas, and register at least one author as a reviewer.
\end{enumerate}

\noindent Before Friday 25 September, AoE: replace \S6 by the protocol; make claims the reported unit, with
39, 15 and 0 as the reason; apply the ten items of Chapter~\ref{ch:review}; replace the admission in \S7
that nothing speaks to discovery error by the estimated-class result; move the clinician sentence to page one; cut
to nine pages; write the AI use statement, which your \texttt{95\_statements.tex} does not have yet.

\begin{sonorio}[Sentences the paper can print]
\begin{itemize}[nosep, leftmargin=*]
\item Counted in the analyst's claims, the certificate did not over-certify on any of 3,014 certifiable
  rows of eleven knowledge conditions on the true CPDAG: whenever the committed set was invalid at the truth, it licensed fewer withdrawn claims than the
  analyst had made false ones. This follows from the definition when the equivalence class and the
  validity criterion are correct, and is not an estimate. Counted in covering hops, the same audit found
  over-certification on 39 of those rows.
\item With one wrong claim, the stated report covered 0.597, and the knowledge interval at one claim
  restored 0.980 at 85\,\% of the calibrated width of discarding the knowledge (23 networks,
  semi-synthetic parameters, $n = 1000$).
\item At matched depth the report differs from a leave-$k$-out screen only where a retraction changes the
  adjustment set while leaving it valid; what it adds is the
  depth the analyst needed and the name of the claim the conclusion rests on.
\item When the class is learned from data, the generating graph extended the analyst's graph on none of 244
  certifiable rows, so the guarantee does not apply.
\item The smallest retraction that breaks the set is a minimal correction set, a notion from the
  satisfiability literature, instantiated here on Meek-closed background knowledge.
\end{itemize}
\end{sonorio}

Each of those sentences has a tempting neighbour that the evidence does not license, and the checklist
below is the list of those neighbours.

\begin{noprato}[Checklist: sentences the paper must not contain]
\begin{itemize}[nosep, leftmargin=*]
\item that the radius, the certificate or the stop rule \emph{beats} a leave-$k$-out screen;
\item that anything is certified robust under an \emph{estimated} CPDAG;
\item a continuity bound («a small knowledge error gives a small bias»);
\item a probability reading of $r$, or that $r$ estimates how wrong an expert is;
\item that zero dangerous rows for the claim certificate is a statistic;
\item that language models are the subject of the paper;
\item the retracted numbers and phrases: «0/791», «$\Ostar$ is a local functional», «up to 128×»;
\item that the MPDAG poset of Taeb, Guo and Henckel is graded.
\end{itemize}
\end{noprato}

\prompt[conceito=frases-proibidas,dificuldade=0.65]{Give four sentences from the checklist that the merged
paper must not contain, and for two of them the measurement that rules them out.}%
{Any four of: the radius or stop rule beats a leave-$k$-out screen; robustness certified under an estimated
CPDAG; a continuity bound; a probability reading of $r$; zero dangerous claim rows presented as a statistic;
language models as the subject; the retracted «0/791», «$\Ostar$ is a local functional», «up to 128×»; the
general MPDAG poset is graded. Measurements: the screen ties at matched depth ($+0.000\ [+0.000, +0.000]$,
M1b); under a learned class the truth extends the analyst's graph on 0 of 244 rows; Proposition 9 of Taeb,
Guo and Henckel proves non-gradedness for $|V| \ge 4$.}

\prompt[conceito=pre-registro,dificuldade=0.6]{What does the ICLR 2027 guideline allow to change about the
authors after the abstract deadline, and until when?}%
{No author can be added or removed after the abstract submission deadline; only the order of authors can
change, and only up to the full paper deadline. So the list, with every author's OpenReview profile, must
be settled by Friday 18 September, AoE.}

\naopasse[decisoes-abertas]{Without looking, write the five steps before Friday in order, and name the one
whose result can no longer be changed after 18 September. If you put the abstract before the framing
decision, reread the section on the three options.}

\subsection*{Exercises}

\begin{exercicio}[metodo=rewriting,conceito=deriva-enquadramento]
Rewrite these two sentences of your abstract for option~A, respecting the checklist and with every number
carrying its population: «On synthetic and on published causal graphs we ask whether the radius ranks
instances by whether the committed set survives corrupted background knowledge, and how the quantity is
distributed. The radius is a diagnostic and not an objective: it prices the structural risk carried by an
adjustment set, in the units the analyst reasons in.»
\end{exercicio}
\begin{solucao}
One possible rewrite. «On 27 networks with background knowledge from elicited and chance-level sources,
audited only after every report was frozen, the report counted in the analyst's claims never
over-certified, while a count of covering hops did so on 39 of 3,014 rows. On 23
networks with semi-synthetic parameters, one wrong claim lowered the coverage of the stated report to 0.60,
and the knowledge interval at one claim restored 0.98 at 85\,\% of the calibrated width of discarding the
knowledge; at matched depth the report takes
the decisions of a leave-$k$-out screen, and it names the claim the conclusion rests on. The guarantee
assumes a correct equivalence class: with classes learned from data, the generating graph extended the
analyst's graph on none of 244 rows.» What changed: «ranks» goes, since ranking is the synthetic panel's
question; «in the units the analyst reasons in» becomes a measured contrast between claims and hops; the
tie is stated rather than hidden; and the class assumption moves into the abstract.
\end{solucao}

\begin{exercicio}[metodo=classification,conceito=frases-proibidas]
Classify each sentence as (L) licensed as written, (C) licensed only with its condition stated, or (F)
forbidden. (a) «The claim certificate never over-certifies.» (b) «The breakdown radius outperforms a
leave-one-out screen.» (c) «With an estimated CPDAG, the committed set remains certified robust.» (d) «One
wrong claim lowered the stated report's coverage from 0.967 to 0.597 (Panel A, $n = 1000$).» (e) «A radius
of 2 means the expert is probably wrong about two orientations.» (f) «The smallest breaking retraction is a
minimal correction set.» (g) «The space of knowledge states is graded.» (h) «Small errors in background
knowledge produce small bias.»
\end{exercicio}
\begin{solucao}
(a) C: true given a correct class and a validity oracle that matches the instrument, and structural rather
than statistical, so both conditions go in the sentence. (b) F: they tie at matched depth, and the
radius-based rule separates only from the depth-one screen. (c) F: the truth extends the analyst's graph on
0 of 244 certifiable rows. (d) L: measured, with its population. (e) F: a probability reading of $r$.
(f) L, provided the citation is there: the name is conceded, not claimed (\texttt{docs/\allowbreak PRIOR\_ART.md}). (g) C: on a
fixed skeleton and under the anti-exchange conjecture; the general MPDAG poset is not graded (Proposition 9).
(h) F: the radius is a breakdown value, not a continuity bound.
\end{solucao}

\section*{Chapter map}
\begin{center}
\begin{tikzpicture}[node distance=0.7cm and 1.7cm]
\node[cmapc] (todo) {your two open\\ experiments};
\node[cmapg, right=of todo] (runs) {ledger and calibration,\\ already run};
\node[cmapr, below=of runs] (conf) {four conflicts};
\node[cmap, below=of todo] (a) {option A: your skeleton,\\ our \S6};
\node[cmap, below=of a] (lock) {18 Sep: authors,\\ abstract, framing};
\node[cmap, below=of conf] (paper) {25 Sep: nine pages,\\ review items};
\node[cmapw, below=0.7cm of $(lock.south)!0.5!(paper.south)$] (sent) {licensed sentences and the checklist};
\draw[cedge] (todo) -- node[above]{closed by} (runs);
\draw[cedge] (runs) -- node[right]{expose} (conf);
\draw[cedge] (conf) -- node[above]{settled in} (a);
\draw[cedge] (a) -- node[left]{first} (lock);
\draw[cedge] (lock) -- node[above]{then} (paper);
\draw[cedge] (lock) -- node[left, pos=0.4]{written with} (sent);
\draw[cedge] (paper) -- node[right, pos=0.4]{checked against} (sent);
\end{tikzpicture}
\end{center}
